% ============================================================
% Delivery Note Template --- VIEWS
% ============================================================
% A delivery note accompanies a set of data handed to someone outside
% the team that produced it. It records WHAT was delivered, HOW to read
% it, and WHAT IS NOT in it. It does not ask the recipient to decide
% anything --- that is a pre-release note's job.
%
% TO USE THIS TEMPLATE:
%   1. Copy this directory to ../<delivery-id>/ and work there.
%   2. Fill the six \delivery* macros below. All six are required.
%   3. Build the \section + \input spine in the DOCUMENT BODY to suit
%      this delivery. There is no fixed structure --- comment in,
%      comment out, reorder, rename. A stub exists for each section in
%      sections/; delete the ones you do not use.
%   4. Write body prose in sections/*.tex. Those files contain NO
%      \section commands --- the headings live here, so reordering is a
%      single-file edit.
%   5. Build with ./build.sh (three pdflatex passes).
%
% Conventions: ../README.md, and the house style guides at
%   ~/brain/2_projects/fao02/_dev_materials/release_notes/STYLE.md
%   ~/brain/2_projects/fao02/_dev_materials/prerelease_notes/STYLE.md
% British spelling. Title Case headings, sentence case bodies.
% booktabs rules only, no vertical rules, every table captioned and
% labelled.
% ============================================================

\documentclass[11pt,a4paper]{article}

% ------------------------------------------------------------
% LANGUAGE & ENCODING
% ------------------------------------------------------------
\usepackage[english]{babel}
\usepackage[utf8]{inputenc}
\usepackage[T1]{fontenc}        % REQUIRED: OT1 cannot set a, o, edh etc.

% ------------------------------------------------------------
% TYPOGRAPHY & SPACING
% ------------------------------------------------------------
\usepackage{microtype}           % Better kerning and justification
\usepackage[onehalfspacing]{setspace} % Line spacing
\usepackage{parskip}             % Paragraph spacing, no indent
\setlength{\footnotesep}{0.7\baselineskip}
\usepackage[hang,flushmargin]{footmisc} % Footnotes: no indent

% ------------------------------------------------------------
% PAGE LAYOUT
% ------------------------------------------------------------
\usepackage[a4paper,margin=1in]{geometry} % Clean margins
\usepackage{fancyhdr}
\pagestyle{fancy}
\fancyhf{} % clear all fields
% Head/foot fields are set after the delivery metadata is defined,
% further down. Inherited \lhead{Maase et al.} / \rhead{\today} removed
% deliberately: the date must not float on rebuild.
\renewcommand{\footrulewidth}{0.8pt}

% ------------------------------------------------------------
% GRAPHICS & TABLES
% ------------------------------------------------------------
\usepackage{graphicx}
\usepackage{booktabs}
\usepackage{tabularx}
\usepackage{xltabular}
\usepackage{array}
% Weighted full-width column: L{w} where the w's across a row sum to the column count.
\newcolumntype{L}[1]{>{\raggedright\arraybackslash\hsize=#1\hsize}X}
\usepackage{makecell}
\usepackage[section]{placeins} % Prevent floats before sections
\usepackage[table]{xcolor}

% ------------------------------------------------------------
% MATH & SYMBOLS
% ------------------------------------------------------------
\usepackage{amsmath,amssymb}

% ------------------------------------------------------------
% CITATION & QUOTING
% ------------------------------------------------------------
\usepackage{csquotes}
\usepackage{natbib}

% ------------------------------------------------------------
% HYPERLINKS
% ------------------------------------------------------------
\usepackage{hyperref}
\hypersetup{colorlinks=true, allcolors=black, urlcolor=[rgb]{0,0.25,0.55}}

% ------------------------------------------------------------
% NOTES
% ------------------------------------------------------------
\usepackage[colorinlistoftodos]{todonotes}

% ------------------------------------------------------------
% WATERMARK
%\usepackage{draftwatermark}
%\SetWatermarkText{DRAFT}
%\SetWatermarkScale{7}          % Size of the text (default: 1)
%\SetWatermarkLightness{0.9}    % 0 = black, 1 = invisible

% ------------------------------------------------------------
% DELIVERY METADATA --- fill all of these
% ------------------------------------------------------------
% IDENTIFIER. Composed from the facts below, never hashed:
%
%     <library>-<loa>-<scope>-<partition>-<YYYY-MM>[-<n>]
%
%   library    the model library that produced it: r2darts2, hydranet,
%              stepshifter, ...
%   loa        the VIEWS level of analysis: pgm (PRIO-GRID month) or
%              cm (country month). This carries BOTH the spatial unit
%              and the temporal resolution.
%   scope      the spatial extent: global, africa, ... Always present.
%              A delivery is not assumed global.
%   partition  calibration / validation / forecasting
%   YYYY-MM    month of delivery
%   n          sequence number, ONLY if a second delivery with the same
%              fields goes out in the same month
%
% CHECKSUM. \deliverychecksum is derived from every specification field
% below by ../checksum.py --- it is not chosen. Regenerate it whenever
% any of those fields changes:
%
%     python3 ../checksum.py . > /dev/null && python3 ../checksum.py .
%
% build.sh verifies it, so a note cannot quietly disagree with its own
% specification.
%
% \deliverydate is deliberately NOT \today. A note records what was sent
% on a given day; rebuilding it must not restamp it. Set it once and
% leave it.
%
% \deliverysha is the views-models commit THE DATA WAS PRODUCED AT ---
% not the commit of this document. The MANIFEST beside each model's
% output records the sha its config was read at; use the latest commit
% that touched any delivered model's configuration. It MUST be pushed,
% or the \cfg links will not resolve.

\newcommand{\deliveryid}{hydranet-pgm-global-calibration-2026-09}
\newcommand{\deliverychecksum}{36172716}
\newcommand{\deliverytitle}{Eight HydraNet Forecast Models, Collapsed Three Ways}
\newcommand{\deliveryauthor}{Simon Polichinel von der Maase}
\newcommand{\deliverydate}{11-10-2026}

% --- SCOPE: what the data is --------------------------------
% Resolution and extent on each axis, stated in parallel. The level-of
% -analysis shorthand (pgm, cm) is not repeated here --- it is the second
% field of the identifier.
\newcommand{\deliverypartition}{Calibration}
\newcommand{\deliveryspatialres}{PRIO-GRID cell}
\newcommand{\deliveryspatialextent}{Global land --- 64,818 cells}
\newcommand{\deliverytemporalres}{Month}
\newcommand{\deliverytemporalextent}{Forecasts Jan 2018 to Dec 2021; trained on Jan 1990 to Dec 2017}
\newcommand{\deliverystructure}{13 rolling origins, each forecasting 36 months, collapsed three ways}
\newcommand{\deliverytargets}{State-based, non-state and one-sided fatalities}

% --- PROVENANCE: what produced it ---------------------------
\newcommand{\deliverylibrary}{views-hydranet, pinned \texttt{\textasciitilde{}=0.1.1}; exact version not recorded}
\newcommand{\deliveryframework}{views-pipeline-core, version not recorded}
\newcommand{\deliverybackend}{views-datafactory, pinned \texttt{>=1.9.0,<2.0.0}; exact version not recorded}
\newcommand{\deliverysha}{1b9cebc}
\newcommand{\deliverycompute}{Rented cloud GPUs (RunPod), NVIDIA RTX PRO 4500}

% Link a model name to its hyperparameter config at the delivery commit.
% Underscores in model names need \catcode`\_=12 around the table that
% uses this --- see the darts note for the pattern.
\newcommand{\cfg}[1]{%
  \href{https://github.com/views-platform/views-models/blob/\deliverysha/models/#1/configs/config_hyperparameters.py}%
       {link}}
\newcommand{\qs}[1]{%
  \href{https://github.com/views-platform/views-models/blob/\deliverysha/models/#1/configs/config_queryset.py}%
       {link}}

% ------------------------------------------------------------
% TITLE
% ------------------------------------------------------------
\usepackage{titling}

\pretitle{\vspace{2em}\centering\Huge\bfseries}
\posttitle{\par\vspace{1em}}
\preauthor{\centering\large}
\postauthor{\par\vspace{0.5em}}
\predate{\centering\small}
\postdate{\par\vspace{2em}}

\title{\deliverytitle}
\author{\deliveryauthor}
\date{\deliverydate}

% ------------------------------------------------------------
% RUNNING HEAD --- the identifier travels on every page, so a loose
% printed page can still be traced back to the delivery it came from.
% ------------------------------------------------------------
\lhead{\deliveryauthor}
\rhead{\deliverydate}
\lfoot{PRIO --- VIEWS \quad \texttt{\deliveryid}}
\rfoot{\thepage}

% ------------------------------------------------------------
% DOCUMENT STARTS
% ------------------------------------------------------------
\begin{document}

% ------------------------------------------------------------
% TITLE PAGE
% ------------------------------------------------------------
\begin{titlepage}
    \centering
    \vspace{3cm}

    {\Huge \textbf{VIEWS}\par}
    {\Huge Delivery Note \par}
    \vspace{0.6cm}
    {\LARGE \deliverytitle \par}
    \vspace{1.5cm}

    \noindent\makebox[\textwidth]{\Large\textemdash\quad$\diamond$\quad\textemdash}
    \vspace{1.5cm}

    \begin{tabular}{@{}r@{\quad}p{10cm}@{}}
        \textbf{Identifier}  & \texttt{\deliveryid} \\
        \textbf{Checksum}    & \texttt{\deliverychecksum} \\[0.35em]
        \textbf{Date}        & \deliverydate \\
        \textbf{Author}      & \deliveryauthor \\[0.9em]
        \multicolumn{2}{@{}l@{}}{\textbf{\textit{Scope}} \quad\small\textit{what the data is}} \\[0.2em]
        \textbf{Partition}            & \deliverypartition \\
        \textbf{Spatial resolution}   & \deliveryspatialres \\
        \textbf{Spatial extent}       & \deliveryspatialextent \\
        \textbf{Temporal resolution}  & \deliverytemporalres \\
        \textbf{Temporal extent}      & \deliverytemporalextent \\
        \textbf{Forecast structure}   & \deliverystructure \\
        \textbf{Targets}              & \deliverytargets \\[0.9em]
        \multicolumn{2}{@{}l@{}}{\textbf{\textit{Provenance}} \quad\small\textit{what produced it}} \\[0.2em]
        \textbf{Model library}        & \deliverylibrary \\
        \textbf{Framework}            & \deliveryframework \\
        \textbf{Data backend}         & \deliverybackend \\
        \textbf{Source commit}        & \texttt{views-models@\deliverysha} \\
        \textbf{Computed on}          & \deliverycompute \\
    \end{tabular}

    \vspace{0.8cm}
    \textit{PRIO -- Oslo}

    \vfill
    \includegraphics[width=0.25\textwidth]{figures/VIEWS_logo_tagline.png}
    \hspace{1cm}
    \includegraphics[width=0.25\textwidth]{figures/PRIO_Logo.png}

    \thispagestyle{empty}
\end{titlepage}

% ============================================================
% DOCUMENT BODY
% ------------------------------------------------------------
% There is NO fixed structure. Comment in what this delivery needs,
% comment out the rest, reorder freely. The headings below are
% suggestions drawn from past deliveries, not a required spine.
% ============================================================

% The spine below is a starting point drawn from past deliveries, not a
% required structure. Comment in, comment out, reorder, rename.
%
% Section 1 is the exception: a delivery note should always open with a
% Context section stating what the document is, what it covers, what it
% explicitly does NOT cover, and that the delivery is a fixed artefact
% identified by its identifier and checksum. That shape is inherited
% from the FAO release notes --- see ../README.md.

\section{Context}
This Delivery Note accompanies \texttt{\deliveryid}, a set of conflict forecasts produced by eight
HydraNet models on the calibration partition at global PRIO-GRID month, each collapsed from its
posterior draws three different ways. It is the authoritative description of what was handed over.
The folder was uploaded on 29 September 2026 with a short README and an audit notebook and no note;
this document supplies the note, and supersedes that README where the two disagree.

The purpose of this note is to:
\begin{itemize}
    \item record what was delivered, in enough detail to identify it again later;
    \item describe the models that produced it, and what distinguishes them from one another;
    \item explain why there are three files per model per origin rather than one, and what each
          one is;
    \item state how the files are organised and how they are intended to be consumed; and
    \item set out what is \emph{not} in the delivery, and what is not yet known about it.
\end{itemize}

\noindent This note covers:
\begin{itemize}
    \item the architecture all eight share, and the eight configurations that differ
          (Sections~\ref{sec:model}--\ref{sec:models});
    \item the three point estimators, the arithmetic of each, and which one the models declare
          (Section~\ref{sec:frames});
    \item the file layout, schema and units of the delivered forecasts
          (Section~\ref{sec:files});
    \item measured magnitude and horizon behaviour --- diagnostics, not an evaluation
          (Section~\ref{sec:magnitude}); and
    \item what the delivery excludes, and the caveats that bear on using it
          (Sections~\ref{sec:excluded}--\ref{sec:caveats}).
\end{itemize}

\noindent This note explicitly does \emph{not} cover:
\begin{itemize}
    \item \textbf{any evaluation against observed outcomes.} Section~\ref{sec:magnitude} reports
          what the models predict and how that varies by forecast step and by estimator. It scores
          nothing, ranks nothing and compares nothing to actuals beyond the single magnitude ratio
          the run itself computed;
    \item any recommendation as to which of the eight, or which of the three estimators, to
          prefer. That is what the selection routine is for;
    \item the validation or forecasting partitions, or any live forecast; and
    \item the internal machinery that produced the run --- the cost, the hardware and what the
          deployment cost us are recorded separately, in
          \texttt{reports/runpod\_cost\_and\_time\_note\_2026-09.md} and
          \texttt{reports/postmortem\_runpod\_first\_deployment\_2026-09.md}.
\end{itemize}

\noindent The delivery is a fixed artefact. It is identified by the identifier and checksum on
the title page --- quote both --- and is not revised in place: any correction is issued as a new
delivery with a new identifier and a new checksum.


\section{The Model}\label{sec:model}
All eight delivered models are the same architecture under eight configurations. The architecture
is \textbf{HydraNet} --- in-house work with no upstream: the design, the implementation and the
paper describing it \citep{maase_hydranet}. The code as it ran is pinned at
\texttt{views-hydranet@e3a039f}; what follows describes it at that commit.

\subsection*{What It Is}

HydraNet treats the global grid as an image sequence rather than as a panel of independent cells.
Each training example is a 32-month window of the whole $360 \times 720$ half-degree grid, with
three input channels --- the state-based, non-state and one-sided fatality histories. The network
is a recurrent U-Net: two encoder levels with batch normalisation and max pooling compress the
grid spatially, four parallel ConvLSTM cells carry temporal memory between months, and skip
connections return the compressed representation to full resolution. One set of parameters is
shared across every cell, so a cell with almost no conflict history is forecast using structure
learned from the whole grid and from its neighbours.

\subsection*{Six Heads, Two Per Target}

The name is the architecture's main feature. There are six independent decoder heads --- two for
each of the three outcomes:

\begin{itemize}
    \item a \textbf{gate} head, trained as a classifier on whether the outcome is non-zero at all,
          emitting $P(y > 0)$ per cell per month; and
    \item a \textbf{body} head, trained as a regression on how large the outcome is, emitting the
          parameters of a count distribution rather than a single number.
\end{itemize}

\noindent The split matters because the overwhelming majority of cell-months carry no fatalities at
all --- in the delivered forecasts, 99.4\% of them (Section~\ref{sec:magnitude}). A single head
asked to produce both the zeros and the magnitudes is pulled toward predicting almost nothing
everywhere; separating the two lets each be trained on the question it is actually answering. The
gate is trained with a weighted binary cross-entropy (positive class weighted $2.0$); the body
with mean squared error on the $\log(1+y)$ scale.

A forecast is the two composed. Two compositions are used across the eight, and the choice is the
single most consequential setting that varies between them:

\begin{itemize}
    \item \textbf{soft gate} --- for each draw, a Bernoulli trial on the gate probability decides
          whether that cell is non-zero in that draw. A cell with a gate of $0.3$ is zero in
          roughly seven draws out of ten and carries the body's magnitude in the other three.
    \item \textbf{threshold gate} --- the cell is kept if the gate exceeds a fixed threshold
          $\tau$ and zeroed otherwise, identically in every draw. Three of the eight use this,
          at $\tau = 0.14$, $0.16$ and $0.20$.
\end{itemize}

\noindent The masks are exactly zero or one and are applied on the $\log(1+y)$ scale, where
$\log(1+0) = 0$, so the multiplication introduces no approximation.

\subsection*{Why There Are Draws At All}

The body head does not emit a number, it emits a distribution --- a negative binomial, or a
mixture of two negative binomials, depending on the model. Drawing from it gives one possible
future; drawing repeatedly gives a sample of them. This delivery carries \textbf{16 draws per cell
per month}, built as four passes with dropout left active at prediction time, each yielding four
draws from the distribution the body emitted ($4 \times 4 = 16$). Both sources of variation are
deliberate: the dropout passes express uncertainty about the model, the family draws express
uncertainty about the outcome given the model.

Sixteen is a setting for this run and not a property of the architecture. The platform standard is
higher; the full-depth configuration peaks at roughly 28.6~GB and does not fit the hardware this
run had. The consequence that bears on the delivered files is in
Section~\ref{sec:frames}: an upper quantile of 16 draws is a coarse statistic.

\subsection*{What Reaches Disk}

The forecasts are produced one month at a time, each month's prediction feeding back as the next
month's input, out to 36 months. Because every model declares \texttt{evaluation\_mode:
stochastic}, the run keeps the full set of draws rather than reducing them, and the collapse to a
single number per cell happens afterwards and outside the run --- which is what
Section~\ref{sec:frames} describes. The draws are composed on the $\log(1+y)$ scale and are on the
count scale by the time they reach the collapse, which refuses them if they are not.


\section{The Eight Models}\label{sec:models}
Nearly everything is held fixed across the eight. All of them are
\texttt{HydraBNUNet06\_LSTM4}; all train for 300 lessons of three windows each at learning rate
$0.001$ with the same warm-up and decay schedule; all take the same three input channels on the
same global land grid of 64,818 cells; all target the same three outcomes on the $\log(1+y)$
scale; all produce 16 draws; all were trained in the same two days on the same rented hardware
from the same commit. Those facts are stated once here and are not repeated per model.

\begingroup
\catcode`\_=12  % underscores are ordinary characters in this table
\begin{table}[h!]
\centering
\small
\renewcommand{\arraystretch}{1.2}
\caption{The eight delivered models, and the five settings that differ between them.}
\label{tab:models}
\begin{tabular}{@{}p{2.8cm}p{1.9cm}p{2.2cm}p{1.6cm}p{2.0cm}p{1.2cm}p{1.1cm}@{}}
\toprule
\textbf{Model} & \textbf{Body} & \textbf{Composition} & \textbf{Window}
  & \textbf{Scheduled} & \textbf{Config} & \textbf{Feat.} \\
 & \textbf{family} & & \textbf{sampling} & \textbf{sampling} & & \\ \midrule
purple_alien     & mixture NB & soft gate        & threshold & off & \cfg{purple_alien}     & \qs{purple_alien} \\
pink_pirate      & mixture NB & soft gate        & threshold & off & \cfg{pink_pirate}      & \qs{pink_pirate} \\
blue_stranger    & mixture NB & soft gate        & Boltzmann & off & \cfg{blue_stranger}    & \qs{blue_stranger} \\
bold_comet       & mixture NB & threshold, 0.14  & threshold & off & \cfg{bold_comet}       & \qs{bold_comet} \\
blazing_meteor   & mixture NB & threshold, 0.16  & threshold & off & \cfg{blazing_meteor}   & \qs{blazing_meteor} \\ \addlinespace
heavy_freighter  & NB         & soft gate        & threshold & \textbf{on} & \cfg{heavy_freighter}  & \qs{heavy_freighter} \\
bright_starship  & NB         & soft gate        & threshold & off & \cfg{bright_starship}  & \qs{bright_starship} \\
violet_visitor   & NB         & threshold, 0.20  & sigmoid   & off & \cfg{violet_visitor}   & \qs{violet_visitor} \\
\bottomrule
\end{tabular}
\end{table}
\endgroup

\noindent The link columns open each model's configuration on GitHub, fixed at the delivery commit
on the title page: \textbf{Config} is \texttt{config\_hyperparameters.py}, which sets the body
family, the composition, the dropout rate and the training schedule; \textbf{Feat.} is
\texttt{config\_queryset.py}, which declares the inputs. The names are arbitrary labels assigned
by the \texttt{adjective\_noun} convention used across views-models and carry no meaning.

Two further settings vary and are not in the table because they are not interpretable as
treatments: the random seed (42--47, locked and logged per run) and the dropout rate (0.125 in six
models, 0.1 in \texttt{blue\_stranger}, 0.15 in \texttt{violet\_visitor}).

\subsection*{The Eight, One Paragraph Each}

\noindent\textbf{\texttt{purple\_alien}} --- Mixture-negative-binomial body, soft gate. The
reference configuration the other seven are variations on, and the only one carrying bespoke
regression-loss weights. It is the model used for the worked examples throughout this note.

\noindent\textbf{\texttt{pink\_pirate}} --- Identical to \texttt{purple\_alien} in body family and
composition, with the standard regression-loss weights and a different seed. The nearest thing in
the set to a seed replicate of the reference configuration.

\noindent\textbf{\texttt{blue\_stranger}} --- Mixture-negative-binomial body, soft gate, and the
only model whose training windows are drawn by Boltzmann weighting: a window's anchor cell is
sampled with probability proportional to $\exp(\text{activity}/\tau)$ rather than uniformly from
cells above an activity threshold. It also carries the lowest dropout rate of the eight, 0.1.

\noindent\textbf{\texttt{bold\_comet}} --- Mixture-negative-binomial body with a threshold gate at
$\tau = 0.14$: a cell is forecast as non-zero only where the gate exceeds that probability, and
identically in all 16 draws. The lowest of the three thresholds, so the most permissive of the
three threshold-gate models.

\noindent\textbf{\texttt{blazing\_meteor}} --- As \texttt{bold\_comet} with the threshold raised to
$\tau = 0.16$. The pair is the closest thing in the set to a clean reading of what moving the
threshold does, though the seeds differ between them.

\noindent\textbf{\texttt{heavy\_freighter}} --- Single-negative-binomial body, soft gate, and
\textbf{the only model trained with scheduled sampling active}, which feeds the model's own
predictions back during training rather than always the observed history. It is retained
deliberately: it was the only configuration predicting fatalities at a realistic scale, and it
remains the highest-magnitude model of the eight (Section~\ref{sec:magnitude}).

\noindent\textbf{\texttt{bright\_starship}} --- Single-negative-binomial body, soft gate,
scheduled sampling off. Intended as the controlled contrast to \texttt{heavy\_freighter}, and the
lowest-magnitude model of the eight. \textbf{The contrast is confounded:} the two also differ in
seed, so the gap between them cannot be attributed to scheduled sampling alone.

\noindent\textbf{\texttt{violet\_visitor}} --- Single-negative-binomial body with the highest
threshold gate, $\tau = 0.20$, the highest dropout rate, 0.15, and the only sigmoid window
sampler, which weights a window's anchor cell by $\sigma(k(\text{activity} - \text{threshold}))$
instead of applying a hard cut.

\subsection*{What the Design Does and Does Not Establish}

The two axes that were chosen as treatments are the body family (mixture of two negative binomials
in five models, a single negative binomial in three) and the composition (soft gate in five,
threshold gate in three). Both are real trades rather than one option dominating: the threshold
gate ranks slightly better on the run's own metrics and predicts measurably less on every model it
was tried on, which is why both are carried into the selection rather than one being chosen here.

Three things about this set are assumptions rather than findings, and are recorded as such: the
five-to-three split between the two body families, the specific seeds, and the premise that eight
separate trainings are more informative than fewer models each carrying more compositions. The
eight are a spread of plausible configurations submitted for selection, not a factorial design
from which an effect can be read.


\section{The Three Frames}\label{sec:frames}
There are three files for every model and every forecast origin. \textbf{They are not three models
and not three runs.} There was one run per model. It produced 16 draws for every cell and every
month, and a point forecast needs one number rather than 16, so something has to collapse them.
There is more than one defensible way to do that, and three were used.

\begin{table}[h!]
\centering
\small
\renewcommand{\arraystretch}{1.3}
\caption{The three frames. One run, one set of draws, collapsed three ways.}
\label{tab:frames}
\begin{tabular}{@{}p{1.1cm}p{3.6cm}p{9.3cm}@{}}
\toprule
\textbf{Frame} & \textbf{Directory} & \textbf{What it is} \\ \midrule
I   & \texttt{I\_mean} &
  The arithmetic mean of the 16 draws, taken in double precision on the count scale.
  \textbf{This is the estimator the models declare}, and the only one of the three that the
  delivery pipeline produces. \\ \addlinespace
II  & \texttt{II\_q95} &
  The 95th percentile across the 16 draws. An upper estimate: the magnitude that is exceeded in
  roughly one draw in twenty. \\ \addlinespace
III & \texttt{III\_conditional\_mean} &
  $\mathrm{E}[y \mid y > 0]$ --- the mean taken over the positive draws only, equivalently
  $\mathrm{E}[y] / P(y>0)$. The expected magnitude \emph{given} that something happens, rather
  than the expected magnitude including the chance that nothing does. \\
\bottomrule
\end{tabular}
\end{table}

\noindent Lexical sorting puts these in the order III, II, I, so a file browser or a Drive listing
will show them backwards. The Roman numerals are in the directory names so that the intended
reading order survives being sorted.

\subsection*{Why the Collapse Happens Outside the Run}

All eight models declare \texttt{evaluation\_mode: stochastic}, which tells the run to keep the
full set of draws and not to reduce them. Each model also declares an aggregation method
(\texttt{arithmetic\_mean}), but under stochastic evaluation that declaration has no effect and the
run says so in its own log. The consequence is that the collapse is a separate, later step, taken
deliberately rather than as a side effect of how the run was configured --- and that is the step
that produced these three frames.

\subsection*{What Frame I Can Be Checked Against, and What Frames II and III Cannot}

This needs stating plainly, because it is the one asymmetry in the delivery.

\textbf{Frame I is reproducible.} It is produced by a collapse step held in views-models, at
\texttt{tools/collapse/}\allowbreak\texttt{collapse\_predictions.py}, which implements the
arithmetic mean and the median and refuses anything else. It takes the mean in double precision, on raw
counts, and refuses the input outright on a non-finite value, a negative value, a disagreement in
draw count between targets, a duplicated cell-month, or a maximum low enough to suggest the values
are still on a log scale. Those refusals are tested.

\textbf{Frames II and III are not reproducible, and no code for them exists.} Neither the
percentile nor the conditional mean is implemented anywhere in views-models; both were computed
once, in an ad-hoc step, and the posterior draws they were computed from are not on any machine we
still have. The draws were written and compressed on the rented GPUs --- the run manifests record
13~GB of raw draws per model compressing to roughly 9~MB --- but they were never transferred off,
and the machines are gone. \textbf{So frames II and III cannot be recomputed, re-derived or
checked against the posterior.} They are what they are: a probe, recorded here as such.

This also corrects a claim made in my own cost memo of 28 September, which stated that the full
posterior had been retained and that a different summary statistic could therefore be produced
without a new run. That is not true on disk. \textbf{Any statistic other than these three needs a
new GPU run.}

\subsection*{Two Properties of These Estimators}

\textbf{Frame II is coarse at 16 draws.} The 95th percentile of 16 values falls between the 15th
and 16th order statistic, so a single draw moves it, and at a sample this small an upper quantile
is biased low. It is reported as computed, with no small-sample correction. If the draw count is
ever raised, frame II must be re-measured before anything relies on it.

\textbf{Frame III inflates the sparse targets most.} Dividing by $P(y>0)$ multiplies a cell that
was positive in one draw of 16 by about 16. The effect is largest where positive draws are rarest,
which is why frame III runs 4.2 times frame I on state-based fatalities but 7.5 times on one-sided
(Section~\ref{sec:magnitude}). One ratio quoted for all three targets would hide that.

\subsection*{How the Three Are Meant to Be Used}

All three go into the selection routine as if they were separate models, and the published
selection criteria decide. That is settled. I am not proposing one of them here, and the frames
are not ranked anywhere in this note.

One provenance correction while the frames are in view: the README uploaded with the folder cites
``ADR-023 Amendment 1'' as the authority for the three frames. No such document exists. ADR-023
governs the collapse and was accepted on 28 September 2026; its only amendment is dated 6 October
2026 and concerns an unrelated library. The authority for frame I is ADR-023 and the commit on the
title page. Frames II and III have no governing document, which is part of why this section says
what it says about them.


\section{The Delivered Files}\label{sec:files}
The delivery is one folder per model, eight in all, each holding three frame directories of
thirteen Parquet files:

\begin{verbatim}
  <model>/
    I_mean/                 predictions_calibration_<ts>_<NN>.parquet  (13 files)
    II_q95/                 "
    III_conditional_mean/   "
\end{verbatim}

\noindent That is $8 \times 3 \times 13 = 312$ files and a little over 180~MB. The timestamp in each
filename is the model's own run, so it is the same within a model and differs between models; the
two-digit suffix \texttt{\_00} to \texttt{\_12} is the forecast origin.

Parquet is a columnar table format. It is not a text format and will not open in a text editor,
but every analysis environment reads it directly, with no conversion step:

\begin{quote}
\small
\textbf{Python} \\
\texttt{import pandas as pd} \\
\texttt{f = "purple\_alien/I\_mean/predictions\_calibration\_20260928\_201214\_00.parquet"} \\
\texttt{df = pd.read\_parquet(f)}

\medskip
\textbf{R} \\
\texttt{library(arrow)} \\
\texttt{f <- "purple\_alien/I\_mean/predictions\_calibration\_20260928\_201214\_00.parquet"} \\
\texttt{df <- read\_parquet(f)}
\end{quote}

\subsection*{Thirteen Origins, and Why They Must Not Be Stacked}

The calibration test window is months 457 to 504 --- January 2018 to December 2021. Each model
forecasts 36 months ahead from each of thirteen successive starting points, one month apart.

\begin{table}[h!]
\centering
\small
\renewcommand{\arraystretch}{1.2}
\caption{The thirteen forecast origins. Each file covers 36 months from its own origin.}
\label{tab:origins}
\begin{tabular}{@{}llll@{}}
\toprule
\textbf{File} & \textbf{Forecast from} & \textbf{Months covered} & \textbf{Calendar} \\ \midrule
\texttt{\_00} & month 456 & 457--492 & Jan 2018 -- Dec 2020 \\
\texttt{\_01} & month 457 & 458--493 & Feb 2018 -- Jan 2021 \\
\texttt{\_02} & month 458 & 459--494 & Mar 2018 -- Feb 2021 \\
\multicolumn{4}{@{}l@{}}{\quad\textit{\ldots one month later at each step \ldots}} \\
\texttt{\_11} & month 467 & 468--503 & Dec 2018 -- Nov 2021 \\
\texttt{\_12} & month 468 & 469--504 & Jan 2019 -- Dec 2021 \\
\bottomrule
\end{tabular}
\end{table}

\noindent \textbf{The windows overlap, so the same cell-month appears in many files.} Months 469
to 492 --- twenty-four of the forty-eight in the span --- appear in \emph{all thirteen} origins,
each time at a different forecast age. Concatenating the thirteen
files and treating the result as a panel produces a table with up to thirteen different values for
the same cell and month, which is not a forecast of anything. Pick one origin, or carry the origin
as an explicit dimension. The audit notebook in the folder prints the overlap as a table if it
helps to see it.

\subsection*{Schema}

Every one of the 312 files holds exactly five columns and \textbf{2,333,448 rows} --- 64,818 cells
$\times$ 36 months, a complete block with no gaps. The schema is byte-identical across all 312.

\begin{table}[h!]
\centering
\small
\renewcommand{\arraystretch}{1.2}
\caption{The five columns, identical in every file.}
\label{tab:columns}
\begin{tabular}{@{}p{4.2cm}p{2.0cm}p{8.0cm}@{}}
\toprule
\textbf{Column} & \textbf{Type} & \textbf{Meaning} \\ \midrule
\texttt{month\_id}           & integer & VIEWS month index; 457 is January 2018 \\
\texttt{priogrid\_id}        & integer & PRIO-GRID cell identifier (see below) \\
\texttt{pred\_lr\_sb\_best}  & double  & predicted state-based fatalities \\
\texttt{pred\_lr\_ns\_best}  & double  & predicted non-state fatalities \\
\texttt{pred\_lr\_os\_best}  & double  & predicted one-sided fatalities \\
\bottomrule
\end{tabular}
\end{table}

\noindent \textbf{\texttt{priogrid\_id} is a cell, not a coordinate.} PRIO-GRID is a $720 \times
360$ grid of half-degree cells covering the globe, numbered from 1 in the south-west corner and
running east before north. To place a row on a map:

\begin{table}[h!]
\centering
\small
\renewcommand{\arraystretch}{1.3}
\caption{Converting \texttt{priogrid\_id} to coordinates. The result is the cell centre.}
\label{tab:priogrid}
\begin{tabular}{@{}p{4.6cm}p{9.4cm}@{}}
\toprule
Row (south to north) & $\texttt{row} = \lfloor (\texttt{priogrid\_id} - 1) / 720 \rfloor$ \\
Column (west to east) & $\texttt{col} = (\texttt{priogrid\_id} - 1) \bmod 720$ \\
Latitude  & $-90 + (\texttt{row} + 0.5) \times 0.5$ \\
Longitude & $-180 + (\texttt{col} + 0.5) \times 0.5$ \\
\bottomrule
\end{tabular}
\end{table}

\noindent In R:

\begin{quote}
\small
\texttt{d\$lat <- -90  + (((d\$priogrid\_id - 1) \%/\% 720) + 0.5) * 0.5} \\
\texttt{d\$lon <- -180 + (((d\$priogrid\_id - 1) \%\%  720) + 0.5) * 0.5}
\end{quote}

\noindent As a check: in \texttt{purple\_alien}'s first origin the largest predicted state-based
value is 58.0 fatalities in month 462 --- June 2018 --- at \texttt{priogrid\_id} 181875, which
resolves to $36.25^\circ$N, $37.25^\circ$E, in northern Syria. The 64,818 cells delivered here are
the land subset of the 259,200 in the full grid; their centres span $55.75^\circ$S to
$83.25^\circ$N.

Values are fatality counts. They are continuous rather than integer, because a mean over draws is
not a whole number, and they are on the original count scale --- not logged, not scaled, not
truncated. Observed values are not included (Section~\ref{sec:excluded}).

\subsection*{Checking the Delivery Without Running Anything}

The folder contains \texttt{audit\_selector\_frames.ipynb}, which audits the folder it sits in. It
opens all 312 files, checks the columns and types, checks for non-finite, negative and duplicated
values, verifies that every file is a complete $64{,}818 \times 36$ block, draws the global maps
and the per-step profiles, and ends in a cell that \emph{raises} rather than prints if anything is
wrong.

\textbf{Its outputs are saved from the run that checked this exact copy of the data}, so it can be
read without executing it. To re-run it, start Jupyter in the folder and Run All; it needs only
\texttt{pandas}, \texttt{pyarrow}, \texttt{numpy} and \texttt{matplotlib}, and reads the model
directories beside it. Nothing else in this delivery needs to be verified by hand, and there is no
checksum file to run.

For the record, rather than as something to take on trust: across all 312 files and all 2.18 billion
delivered values there is \textbf{no missing, negative or infinite value}, every file is a complete
$64{,}818 \times 36$ block with no repeated \texttt{(month\_id, priogrid\_id)} pair, and the
columns and types are exact in every file. Two further checks were run for this note. The three
frames are written in the same row order, so they can be compared row by row without a join. And on
every one of the 242,678,592 rows, in all three targets, frame II and frame III are both greater
than or equal to frame I --- no exceptions, worst violation exactly zero. The second of those is
the one worth having: since $\text{mean} = P(y>0) \times \mathrm{E}[y \mid y>0]$, a row where the
conditional mean fell below the mean would mean the two frames had been collapsed from different
draws. None did.

One gap to record: the delivered folder carries \textbf{no per-model manifest.} The engine version,
run host, runtime and configuration commit for each model exist, in the run records held in
views-models, and the commit is on the title page of this note --- but they do not travel inside
the folder. A later delivery of this kind should ship them beside the data.


\section{Magnitude, Agreement and Horizon}\label{sec:magnitude}
\textbf{This section ranks nothing and recommends nothing.} It reports what the delivered files
contain and how that varies by model, by estimator and by forecast step. One quantity here is
measured against observed outcomes --- the magnitude calibration ratio, which the run itself
computed. Everything else is a property of the predictions alone. It is here because these three
behaviours determine whether a given file is being read for the thing it can support, and none of
them is visible from the files without computing it.

\subsection*{Magnitude: All Eight Under-Predict, By Factors That Differ Fourfold}

The magnitude calibration ratio is $\mathrm{MCR} = \bar{\hat{y}} / \bar{y}$, the mean prediction
over the mean observed value on the same cell-months. One is perfect aggregate calibration, below
one is systematic under-prediction. Because it averages over the draws before dividing, it is
exactly the ratio frame I's values bear to the observed totals.

\begingroup
\catcode`\_=12
\begin{table}[h!]
\centering
\small
\renewcommand{\arraystretch}{1.2}
\caption{Magnitude calibration and CRPS on the calibration test window, as computed by the run.
Ordered by state-based MCR. Lower CRPS is better; MCR nearer 1 is better calibrated.}
\label{tab:magnitude}
\begin{tabular}{@{}p{3.3cm}rrrcr@{}}
\toprule
\textbf{Model} & \textbf{MCR sb} & \textbf{MCR ns} & \textbf{MCR os} & & \textbf{CRPS sb} \\ \midrule
bright_starship  & 0.067 & 0.057 & 0.059 & & 0.1143 \\
blue_stranger    & 0.073 & 0.063 & 0.040 & & 0.1136 \\
purple_alien     & 0.093 & 0.127 & 0.095 & & 0.1130 \\
violet_visitor   & 0.110 & 0.169 & 0.106 & & 0.1132 \\
pink_pirate      & 0.195 & 0.135 & 0.060 & & 0.1115 \\
heavy_freighter  & 0.215 & 0.303 & 0.241 & & 0.1135 \\
bold_comet       & 0.240 & 0.171 & 0.257 & & 0.1107 \\
blazing_meteor   & 0.250 & 0.319 & 0.328 & & 0.1101 \\ \midrule
\textit{range}   & \textit{3.8$\times$} & \textit{5.6$\times$} & \textit{8.2$\times$} & & \textit{3.8\%} \\
\bottomrule
\end{tabular}
\end{table}
\endgroup

\noindent Two things in that table are worth stating rather than leaving to be noticed.

\textbf{Every model under-predicts magnitude, and frame I under-predicts it by between four and
fifteen times.} The least timid of the eight reaches a quarter of the observed total; the most
timid reaches a fifteenth. The shortfall is in severity rather than in location: the models mark
roughly the right places and put numbers that are too small in them. This is a property of the
estimator and the objective, not of this export --- a body trained on mean squared error over a
field that is 99.4\% zeros is rewarded for staying near zero.

\textbf{CRPS does not separate the eight the way magnitude does.} Across the eight it spans 3.8\%
while state-based magnitude spans 3.8-fold. The two orderings also run opposite ways here: the most
timid model carries the worst CRPS and the least timid the best. A 3.8\% spread is a narrow basis
for a decision, and the selection criteria --- not this note --- determine what weight it carries.

\subsection*{The Three Estimators, Per Target}

Totals over all eight models and all thirteen origins, with each frame's ratio to frame I.

\begin{table}[h!]
\centering
\small
\renewcommand{\arraystretch}{1.2}
\caption{Total predicted fatalities per frame and target, summed over 8 models $\times$ 13 origins,
with the ratio to frame I.}
\label{tab:frameratios}
\begin{tabular}{@{}lrrrrr@{}}
\toprule
 & \multicolumn{1}{c}{\textbf{I: mean}} & \multicolumn{2}{c}{\textbf{II: q95}}
 & \multicolumn{2}{c}{\textbf{III: E[y|y>0]}} \\
\cmidrule(lr){2-2}\cmidrule(lr){3-4}\cmidrule(lr){5-6}
\textbf{Target} & total & total & ratio & total & ratio \\ \midrule
State-based & 4,194,894 & 18,898,471 & 4.51$\times$ & 17,815,285 & 4.25$\times$ \\
Non-state   &   731,151 &  3,522,070 & 4.82$\times$ &  4,417,776 & 6.04$\times$ \\
One-sided   &   427,970 &  1,979,497 & 4.63$\times$ &  3,217,289 & \textbf{7.52}$\times$ \\ \midrule
All three   & 5,354,015 & 24,400,038 & 4.56$\times$ & 25,450,349 & 4.75$\times$ \\
\bottomrule
\end{tabular}
\end{table}

\noindent Frames II and III both land near five times frame I in total, which is close to the
factor by which frame I falls short of observed --- that is what makes them interesting and is the
reason all three were kept. But the per-target columns are the point of the table: frame III runs
4.25 times frame I on state-based fatalities and 7.52 times on one-sided, for the reason given in
Section~\ref{sec:frames}. A single ratio quoted for all three targets would be wrong for two of
them.

\textbf{And the two estimators are not equally well summarised by a ratio.} Taken per model and
per target --- twenty-four figures --- frame II sits between 4.25 and 4.91 times frame I, so
``about $4.6\times$'' describes every one of them. Frame III ranges from \textbf{2.88 to 11.62}
over the same twenty-four: 2.88 to 6.89 on state-based, 4.19 to 10.71 on non-state, 4.75 to 11.62
on one-sided. Its pooled 4.75 is an average over a fourfold spread and should not be used as a
conversion factor for any particular model.

\textbf{All three frames are non-zero on exactly the same rows.} 1,469,004 of the 242,678,592 rows
in each frame carry a non-zero prediction --- 0.61\%, so 99.4\% of the delivered field is exact
zeros. Those rows fall on 98,995 distinct cell-months and 6,552 distinct PRIO-GRID cells, a tenth
of the 64,818 delivered. The three estimators differ in magnitude and never in where conflict is
predicted; this was verified across all 312 files, not sampled. It is a property of the
construction rather than a coincidence: all three estimators collapse the \emph{same} draw set, and
a cell with no positive draws among its sixteen gives a mean, a 95th percentile and a conditional
mean that are all zero. Note that the figure 1,469,004 counts \emph{rows}, which means the same
cell-month is counted once per model and once per origin; it is not a count of distinct
cell-months.

\subsection*{Horizon: The Field Halves Over Five Steps, Then Stops Changing}

Total predicted state-based fatalities by forecast step, averaged over the thirteen origins and
summed over the eight models.

\begin{table}[h!]
\centering
\small
\renewcommand{\arraystretch}{1.2}
\caption{Predicted state-based fatalities by forecast step. Frame I's step-1 value is 100\%.}
\label{tab:horizon}
\begin{tabular}{@{}rrrrr@{}}
\toprule
\textbf{Step} & \textbf{I: mean} & \textbf{II: q95} & \textbf{III: E[y|y>0]}
 & \textbf{I as \% of step 1} \\ \midrule
 1 & 16,404 & 71,160 & 56,415 & 100.0\% \\
 2 & 11,883 & 52,554 & 46,616 &  72.4\% \\
 3 & 10,058 & 44,794 & 41,453 &  61.3\% \\
 4 &  9,455 & 42,590 & 39,164 &  57.6\% \\
 5 &  8,962 & 40,480 & 38,155 &  54.6\% \\
 6 &  8,789 & 39,563 & 37,213 &  53.6\% \\ \addlinespace
 9 &  8,468 & 38,259 & 36,754 &  51.6\% \\
12 &  8,393 & 37,860 & 35,606 &  51.2\% \\
18 &  8,674 & 39,266 & 37,939 &  52.9\% \\
24 &  8,607 & 38,698 & 36,794 &  52.5\% \\
30 &  8,480 & 38,353 & 36,287 &  51.7\% \\
36 &  8,728 & 39,592 & 37,336 &  53.2\% \\
\bottomrule
\end{tabular}
\end{table}

\noindent \textbf{The total settles after about five steps.} Between step 1 and step 5 the
predicted total falls by 45\%. From step 6 to step 36 --- thirty-one of the thirty-six months
delivered --- frame I stays between 8,265 and 8,907, a spread of 7.8\% with no trend, and the other
two frames behave the same way.

\textbf{What settles is the total, not the field.} These are different claims and only the first
one is true, so the distinction is worth making explicitly. Measuring how much of the field moves
between consecutive forecast steps ---
$\sum_i |f_i(h) - f_i(h-1)| \,/\, \sum_i f_i(h-1)$, the share of the predicted mass that
relocates --- gives a median of \textbf{0.93} pooled over the eight models and all thirteen
origins, and \textbf{0.94} over steps 12 to 36 alone. A frozen field would read 0. Roughly as much
of the field relocates between month 35 and month 36 as between month 1 and month 2. Per model the
figure runs from 0.67 to 1.32, and it rises slightly with the horizon rather than falling.

So the forecasts do not go flat. What goes flat is their marginal intensity: after five steps the
models predict about the same amount of conflict in total, and about the same number of active
cells, while continuing to move it around the map at full amplitude.

\textbf{What this does not establish, and it is the part that matters.} A field that keeps moving
is equally consistent with real month-to-month dynamics and with churn around a fixed
climatological intensity, and this measurement cannot separate the two. It is a property of the
predictions with no observed values in it. So nothing here licenses either reading --- neither that
month 30 carries forecast information nor that it does not. The reason to be careful with the long
horizon is Section~\ref{sec:excluded}: \textbf{nothing in this delivery has been scored at any
horizon}, and a 36-month forecast whose accuracy has not been measured past month one is untested
rather than disproven. A metric averaged evenly over the 36 months is still worth treating with
care, because 31 of those months sit in the regime where the total has stopped changing.

One further thing this does not establish. An earlier scrutiny of these eight, on a smaller region
and a shorter training budget, found per-step magnitude decaying toward zero rather than settling.
That was a different run; this measurement is of these files. The two are different quantities and
neither should be read into the other.


\section{Not Included in This Delivery}\label{sec:excluded}
\begin{table}[h!]
\centering
\small
\renewcommand{\arraystretch}{1.3}
\caption{What this delivery does not contain, and where it is instead.}
\label{tab:not-included}
\begin{tabular}{@{}p{4.4cm}p{10.2cm}@{}}
\toprule
\textbf{Not included} & \textbf{Where it is, or why not} \\ \midrule
Observed values &
  The files carry predictions only. Actuals come from the data store, keyed on the same
  \texttt{month\_id} and \texttt{priogrid\_id}, so a join needs nothing from here. \\ \addlinespace
The posterior draws &
  Produced and compressed on the rented machines, never transferred off them, and the machines are
  gone. This is the one omission with consequences --- Section~\ref{sec:frames}. \\ \addlinespace
Evaluation metrics &
  Beyond the magnitude and CRPS figures in Section~\ref{sec:magnitude}, which the run computed on
  its own output, nothing has been scored. The per-step and per-origin evaluation outputs were
  written on the machines and did not come home either. \\ \addlinespace
Per-model manifests &
  No provenance file travels inside the folder. Run host, runtime, engine version and
  configuration commit are held in views-models; the commit is on the title page. \\ \addlinespace
Training logs &
  Held in views-models beside each model's output, not shipped. They carry the metric tables
  Section~\ref{sec:magnitude} quotes. \\ \addlinespace
Other partitions &
  Calibration only. No validation partition, no forecasting partition, no live forecast. \\ \addlinespace
Model artefacts &
  No trained weights. None from this run exist on any machine we hold, so the run cannot be
  resumed or re-predicted from --- only repeated from the configuration. \\
\bottomrule
\end{tabular}
\end{table}


\section{Caveats and Notes}\label{sec:caveats}
\textbf{Nothing here has been evaluated in the sense the selection routine means.} The figures in
Section~\ref{sec:magnitude} are the ones the run computed on its own output, plus three
measurements taken from the delivered files for this note. No model has been scored against the
published criteria, ranked against another, or compared with any model outside these eight.

\textbf{Two of the three frames cannot be reproduced.} Frame I is produced by code in views-models
and its refusals are tested. Frames II and III were computed once, outside that code, from draws
that were never transferred off the rented machines. They cannot be recomputed or checked, and no
governing document describes them. Section~\ref{sec:frames} sets this out in full; it is repeated
here because it is the caveat most likely to matter later.

\textbf{A statistic other than these three needs a new run.} My cost memo of 28 September said the
full posterior had come home and that a different summary form could be produced without paying
for more compute. That was wrong: the draws are not on disk. Anything else --- a different
quantile, a different conditional, a predictive interval --- requires training to be repeated on
rented hardware.

\textbf{The software is pinned as ranges, not versions.} The configuration commit on the title page
is exact, and the run logs record the CUDA build. The engine versions are not recorded anywhere:
the models pinned \texttt{views-hydranet~}$\sim$\texttt{=0.1.1} and
\texttt{views-datafactory>=1.9.0,<2.0.0}, the environment that resolved those pins existed only on
the rented machines, and no version line was captured before they were destroyed. So the code is
identified to a repository state but not to a release. A later delivery should record the resolved
versions in a manifest beside the data.

\textbf{Pinning the code does not pin the input.} The predictors were fetched live from the data
store at run time. A re-run from the same commit against a later state of that store can produce
different numbers without anything in the configuration having changed. What is fixed is this
delivery, by its identifier and checksum.

\textbf{The horizon behaviour has a mechanism, and it is the gate.} The gate probability decays
with forecast step --- roughly thirtyfold by month 36, in all eight models --- and that decay is
what flattens the field after about five steps. It is not caused by the choice of regression loss;
substituting the loss was tried and changed the decay very little. The predicted magnitude
conditional on firing is comparatively stable; what fades is the model's willingness to fire at
all.

\textbf{One comparison in the set looks controlled and is not.} The models
\texttt{heavy\_freighter} and \texttt{bright\_starship} share a body family and a composition
and differ in scheduled sampling, which invites reading the gap between them as the effect of
scheduled sampling. They also differ in random seed. With one run each, the two explanations
cannot be separated.

\textbf{An earlier scrutiny of these eight reached a harsher verdict, on a different run.} In
August the same eight were examined at a shorter training budget on Africa and the Middle East
only, and that report recommended against promotion on the grounds that seven of the eight
collapsed to near-zero magnitude within six months. This delivery is a different run --- full
training budget, global land --- and Section~\ref{sec:magnitude} measures the behaviour afresh on
the files actually delivered. The August report is superseded for these files and is not
superseded as a record of what it measured.

\textbf{The model carries no covariates.} Each model sees the three fatality histories and nothing
else --- no population, no terrain, no governance indices, no event data beyond the outcomes
themselves. These are purely autoregressive spatiotemporal models. Whatever structure they use is
learned from the spatial and temporal pattern of past fatalities.


\section{Next Steps}\label{sec:next}
Four things follow from this delivery. None of them needs a decision from the reader.

\begin{enumerate}
    \item \textbf{All twenty-four go into the selection routine as separate models.} Eight models
          times three frames. The published criteria decide which collapse wins, rather than my
          choosing one here and then defending it. This was settled in September and is restated
          rather than reopened.

    \item \textbf{The column names in this delivery differ from the darts delivery's, and the
          consuming tool must be told which it is reading.} Here the columns are
          \texttt{pred\_lr\_sb\_best}, \texttt{pred\_lr\_ns\_best} and
          \texttt{pred\_lr\_os\_best}; the darts delivery uses \texttt{pred\_lr\_ged\_sb} and its
          siblings. \texttt{lr\_ged\_\{sb,ns,os\}} is the canonical form and future deliveries
          will carry it, renamed at send-off rather than in the files on disk. Until then the
          target column has to be set explicitly wherever these files are read, not inferred from
          the schema --- a tool that guesses will silently read the wrong delivery's convention.

    \item \textbf{The next delivery of this kind ships its provenance inside the folder.} A
          manifest per model --- resolved engine versions, run host, runtime, configuration
          commit --- and the per-step evaluation outputs retained rather than left on the machine
          that produced them. Both gaps in this delivery trace to the same cause, which is that
          nothing was brought down from the rented hardware except the predictions.

    \item \textbf{The posterior draws are retained on any future run before anything is derived
          from them.} The reason frames II and III are unverifiable is that the draws were
          compressed on a machine and then lost with it. The compression ratio is large enough ---
          roughly 9~MB per model --- that there was never a cost argument for leaving them
          behind.
\end{enumerate}

\noindent Questions about the data, the models or anything in this note come to me.


% ------------------------------------------------------------
% BIBLIOGRAPHY --- build.sh runs bibtex when this is uncommented
% and something is actually cited.
% ------------------------------------------------------------
\bibliographystyle{apalike}
\bibliography{references}

\end{document}
